% ID: SEL-D0010-BA-E02
% ESTADO: cotejada. Segunda pasada visual completa y borrador compilado revisados.
% Fuente: Selectividad/matematicas_aplicadas/EXTRA_TITULAR_MAT-APLIC-CCSS-JULIO-IMPRESO.pdf
% SHA-256: 548f60ee334516c7fa918dba3c28fb3ea5fe753cf76eb05402d7b8acb5a3b560
% Página: 1; bloque A; ejercicio original 2.

Se consideran las siguientes inecuaciones:
\[
5x-4y\leq -19\qquad 3x-4y\leq -13\qquad x\geq -7\qquad -x-y\geq 2
\]
\begin{enumerate}[label=\alph*)]
  \item \textbf{(1.5 puntos)} Represente la región factible definida por las
  inecuaciones anteriores y determine sus vértices.
  \item \textbf{(0.5 puntos)} \textquestiondown{}Cuáles son los puntos en los
  que se alcanzan el mínimo y el máximo de la función
  $G(x,y)=-\dfrac{1}{5}x+\dfrac{5}{2}y$ en la citada región factible?
  \textquestiondown{}Cuál es su valor?
  \item \textbf{(0.5 puntos)} Responda de forma razonada si la función
  $G(x,y)=-\dfrac{1}{5}x+\dfrac{5}{2}y$ puede alcanzar el valor $\dfrac{47}{3}$
  en la región factible hallada.
\end{enumerate}
